\chapter{Your Toolkit and Environment Setup}
\label{ch:toolkit}

\chapterguide%
  {Create folders and run basic \textbf{terminal commands}.}
  {organize data and notebooks, configure an environment, install dependencies, check that everything works, and run labs reproducibly.}

Every lab in this course runs from the same folder, in the same Python environment. Twenty minutes spent here saves hours of chasing path errors later.

\section{The course project layout}

Your project root is the course folder. In this repository it is called \texttt{ARTIFICIAL INTELLIGENCE/}; name your own copy without spaces, for example \texttt{UTP-ARTIFICIAL-INTELLIGENCE/}. It holds these folders:

\noindent\textbf{Datasets}
\begin{lstlisting}[style=utpterminal]
all_datasets/
  lab01-03_airports.csv
  lab01-03_routes.csv
  lab04-AirQuality.csv
  lab05-sensor_readings_24.csv
  lab06-ocr-handwritten-prescription.zip
  lab07-Tweets.csv
  lab08-insurance.csv
\end{lstlisting}

\noindent\textbf{Experiments, helper code, and generated evidence}
\begin{lstlisting}[style=utpterminal]
all_experiments/
  lab01_search_foundations_worked.ipynb        one notebook per lab, each with a .py twin
  lab02_informed_search_worked.ipynb
  lab03_search_engineering_worked.ipynb
  lab04_fuzzy_anfis_worked.ipynb
  lab05_robotics_sensor_control_worked.ipynb
  lab06_handwriting_ocr_worked.ipynb
  lab07_airline_sentiment_worked.ipynb
  lab08_genetic_algorithm_regression_worked.ipynb
  demo_robot_maze.py                           the maze of Chapters 2-4
  demo_sliding_puzzle.py                       the 8-puzzle of Chapters 2-3
  demo_local_search.py                         eight queens and beam search, Chapters 9-10
  demo_game_search.py                          noughts and crosses, Chapter 11
  utils/
    core.py
    openflights.py
    plotting.py
    __main__.py
    test_search.py
    search_queries.csv
results/
  lab01/ ... lab08/
requirements.txt
\end{lstlisting}

\begin{keyidea}
The notebooks import reusable search code from \texttt{utils/}. Each notebook writes its evidence to \texttt{results/}, which should not be edited by hand; \texttt{requirements.txt} lists the packages installed in \texttt{.venv/}. The \texttt{demo\_} scripts need no dataset and print every number the chapters quote for the maze, puzzle and game examples.
\end{keyidea}

\section{The datasets}

Check \texttt{all\_datasets/} first. If your copy of the course already holds the seven files listed above, skip to the next section. Otherwise click a \textbf{dataset name} below to download it, extract the archive, and save the file under the name in the third column. The notebooks look for exactly those names.

\begin{utptable}
\footnotesize
\renewcommand{\arraystretch}{1.1}
\begin{tabularx}{\linewidth}{@{}m{0.36\linewidth}>{\centering\arraybackslash}m{0.07\linewidth}X@{}}
\toprule
\rowcolor{UTPPaleBlue!92}
\utpth{Kaggle dataset} & \utpth{Lab} & \utpth{Save in \texttt{all\_datasets/} as} \\
\midrule
\textbf{OpenFlights:}
\href{https://www.kaggle.com/datasets/open-flights/flight-route-database}{\textbf{Routes}}
\& \href{https://www.kaggle.com/datasets/open-flights/airports-train-stations-and-ferry-terminals/data}{\textbf{Airports}}
& 1--3 & \texttt{lab01-03\_routes.csv}, \texttt{lab01-03\_airports.csv} \\
\href{https://www.kaggle.com/datasets/fedesoriano/air-quality-data-set}
{\textbf{Air Quality}} & 4 & \texttt{lab04-AirQuality.csv} \\
\href{https://www.kaggle.com/datasets/uciml/wall-following-robot}
{\textbf{Wall-Following Robot}} (24 sensors) & 5 & \texttt{lab05-sensor\_readings\_24.csv} \\
\href{https://www.kaggle.com/datasets/mrdude20/doctor-handwriting-recognition-dataset}
{\textbf{Doctor Handwriting}} & 6 & \texttt{lab06-ocr-handwritten-prescription.zip}, kept zipped \\
\href{https://www.kaggle.com/datasets/chandana890/twitter-airline-sentiment-dataset}
{\textbf{Twitter Airline Sentiment}} & 7 & \texttt{lab07-Tweets.csv} \\
\href{https://www.kaggle.com/datasets/mirichoi0218/insurance}
{\textbf{Medical Cost}} & 8 & \texttt{lab08-insurance.csv} \\
\bottomrule
\end{tabularx}
\end{utptable}

\begin{pitfall}
For \textbf{\texttt{FileNotFoundError}}, check the working directory and the filename, including capitals. Fix the path or the name rather than moving files into the project root.
\end{pitfall}

\section{Two ways to set up: plain Python or Anaconda}

\begin{analogy}
A \term{virtual environment} is a separate toolbox for one project. Install a new wrench in this project's box and every other project's box stays exactly as it was, so a package you upgrade for one course cannot quietly break another.
\end{analogy}

\begin{keyidea}
Use one environment per project, and activate it before installing or running packages. Follow Route~A or Route~B below, not both.
\end{keyidea}

\section{Route A: install Python and configure PATH}

Download Python from the \href{https://www.python.org/downloads/}{\textbf{official downloads page}}.

\begin{description}[leftmargin=1.15in,style=nextline]
  \item[Windows] Run the installer, enable \textbf{\texttt{PATH}}, and reopen the terminal. Use \textbf{\texttt{py}}.
  \item[macOS] Run the \textbf{\texttt{.pkg} installer} and restart Terminal. Keep the \textbf{system Python} unchanged.
  \item[Linux] Install \textbf{Python 3} with your package manager. Do not replace the \textbf{system Python}.
\end{description}

\textbf{Windows PATH fix.} If \texttt{py} is missing, open \textbf{Edit environment variables for your account}, then \textbf{Path > Edit}. Add the installer's command directory, for example:

\begin{lstlisting}[style=utpterminal]
%LocalAppData%\Python\bin
\end{lstlisting}

\textbf{macOS PATH fix.} If the installer placed Python in \texttt{/usr/local/bin} but your terminal cannot find it, add that directory to your Z shell's \textbf{PATH}:

\begin{lstlisting}[style=utpterminal]
echo 'export PATH="/usr/local/bin:$PATH"' >> ~/.zshrc
source ~/.zshrc
\end{lstlisting}

For \textbf{Debian or Ubuntu}:

\begin{lstlisting}[style=utpterminal]
sudo apt update
sudo apt install python3 python3-venv python3-pip
\end{lstlisting}

\textbf{Verify} the installation and executable location:

\begin{lstlisting}[style=utpterminal]
# Windows
py --version
py -m pip --version
where py

# macOS or Linux
python3 --version
python3 -m pip --version
which -a python3
\end{lstlisting}

\section{Route A: create an isolated virtual environment}

\textbf{\texttt{venv}} keeps this project's packages separate from other projects. See the \href{https://docs.python.org/3/library/venv.html}{\textbf{venv documentation}}.

The root-level \textbf{\texttt{requirements.txt}} lists the dependencies. A CPU is sufficient for every lab, and no lab needs a GPU.

\begin{lstlisting}[style=utpterminal]
numpy
pandas
matplotlib
scipy
scikit-learn
pillow
pytesseract
pytest
jupyter
\end{lstlisting}

\textbf{Create} the environment in the \textbf{project root}:

\begin{lstlisting}[style=utpterminal]
# Windows
py -m venv .venv

# macOS or Linux
python3 -m venv .venv
\end{lstlisting}

\textbf{Activate} it from the same folder. Your prompt then starts with \texttt{(.venv)}:

\begin{lstlisting}[style=utpterminal]
# Windows (PowerShell)
.venv\Scripts\Activate.ps1

# macOS/Linux
source .venv/bin/activate
\end{lstlisting}

\textbf{Install dependencies} after activation:

\begin{lstlisting}[style=utpterminal]
# Windows
py -m pip install --upgrade pip
py -m pip install -r requirements.txt

# macOS or Linux
python3 -m pip install --upgrade pip
python3 -m pip install -r requirements.txt
\end{lstlisting}

\textbf{Install once; activate every session.} Reuse \textbf{\texttt{.venv}} without reinstalling unchanged requirements.

\begin{utpfig}[The environment routine]{The environment routine. Creating \texttt{.venv} and installing the requirements happen once; every later session starts by activating the environment and goes straight back to the notebooks.}{fig:venv-routine}
\begin{tikzpicture}[
  font=\small,
  pystep/.style={draw=UTPBlue, fill=white, rounded corners=2pt, align=center,
               inner sep=4pt, minimum width=3.3cm, minimum height=1.05cm,
               font=\scriptsize\ttfamily\color{UTPNavy}},
  pyonce/.style={pystep, fill=UTPPaleBlue},
  pyevery/.style={pystep, fill=UTPPaleTeal, draw=UTPTeal},
  pycap/.style={font=\scriptsize\bfseries\color{UTPMuted}},
  pynote/.style={font=\scriptsize\color{UTPMuted}, align=center},
  pygo/.style={-{Stealth[length=2.2mm]}, thick, draw=UTPTeal}]

\node[pyonce]  (venv) at (0,0)     {py -m venv .venv};
\node[pyevery] (act)  at (4.05,0)  {activate};
\node[pyonce]  (inst) at (8.10,0)  {pip install -r\\requirements.txt};
\node[pyevery] (work) at (12.15,0) {run notebooks in\\all\_experiments/};

\draw[pygo] (venv) -- (act);
\draw[pygo] (act)  -- (inst);
\draw[pygo] (inst) -- (work);

\node[pycap] at (0,0.92)     {once per project};
\node[pycap] at (4.05,0.92)  {every session};
\node[pycap] at (8.10,0.92)  {once};
\node[pycap] at (12.15,0.92) {every session};

% the return path: a later session rejoins at activation, not at install
\draw[pygo, draw=UTPOrange]
  (work.south) -- ++(0,-0.85) -| (act.south);
\node[pynote, text=UTPOrange] at (8.10,-1.15)
  {next session: \textbf{reactivate}, do not reinstall};

\end{tikzpicture}
\end{utpfig}

\begin{keyidea}
\textbf{Activate first, then install.} Keep course packages inside the \textbf{project environment}, not in a shared global installation.
\end{keyidea}

\begin{utpfig}[One environment per project]{Separate environments, separate packages. Each project's \texttt{.venv} holds its own copies of the libraries, and the system Python is left alone; activate the intended project environment before running \texttt{pip install}.}{fig:venv-separate}
\begin{tikzpicture}[
  font=\small,
  envbox/.style={draw=UTPBlue, fill=UTPPaleBlue, rounded corners=2pt,
                 align=center, inner sep=4pt, minimum width=4.3cm},
  sysbox/.style={draw=UTPMuted, fill=UTPPaleGray, rounded corners=2pt,
                 align=center, inner sep=4pt, minimum width=4.3cm},
  pkg/.style={font=\scriptsize\ttfamily, align=center},
  lbl/.style={font=\scriptsize\color{UTPMuted}, align=center},
  hit/.style={-{Stealth[length=2.2mm]}, thick, draw=UTPTeal}]

\node[sysbox] (sys) at (0,2.45) {system Python};
\node[pkg,text=UTPMuted,anchor=west] at (2.35,2.45) {leave this alone};

\node[envbox] (a) at (-2.80,0) {\texttt{this course/.venv}};
\node[pkg,anchor=north] at (-2.80,-0.52) {numpy\\pandas\\scikit-learn\\jupyter};

\node[envbox] (b) at (2.80,0) {\texttt{some other project/.venv}};
\node[pkg,anchor=north] at (2.80,-0.52) {its own versions\\of the same names};

\draw[hit] (sys) -- (a);
\draw[hit] (sys) -- (b);

\end{tikzpicture}
\end{utpfig}

\section{Route B: Anaconda or Miniconda}

\textbf{Miniconda} supplies Python and \texttt{conda}; \textbf{Anaconda} adds Navigator and preinstalled packages. Download \href{https://docs.anaconda.com/miniconda/}{\textbf{Miniconda}} or \href{https://www.anaconda.com/download}{\textbf{Anaconda}}.

\begin{description}[leftmargin=1.15in,style=nextline]
  \item[Windows] Use default installer options without adding conda to PATH. Open \textbf{Anaconda Prompt} from Start.
  \item[macOS / Linux] Run the \texttt{.pkg} or \texttt{.sh} installer, then close and reopen the terminal. The prompt shows \texttt{(base)}, which means conda is active.
\end{description}

\textbf{Create} one environment for the course, with a Python version the libraries support, then \textbf{activate} it:

\begin{lstlisting}[style=utpterminal]
conda create -n ai python=3.12
conda activate ai
\end{lstlisting}

\textbf{Install dependencies} with the same \texttt{requirements.txt} used in Route A:

\begin{lstlisting}[style=utpterminal]
pip install -r requirements.txt
\end{lstlisting}

Create once, activate each session, and avoid installing into \texttt{(base)}. Check the active environment and Python:

\begin{lstlisting}[style=utpterminal]
conda env list
python -c "import sys; print(sys.executable)"
\end{lstlisting}

\begin{pitfall}
Avoid installing the same package with both conda and pip. Use pip with the course requirements. Recreate a broken environment with \texttt{conda env remove -n ai}, then repeat setup.
\end{pitfall}

\section{Check that everything works}

Run these three commands from the course folder with your environment active. Each one tests something different, so a failure tells you where to look.

\begin{lstlisting}[style=utpterminal]
python -c "import numpy, pandas, scipy, sklearn, matplotlib, PIL, pytest; print('Packages OK')"
python all_experiments/demo_robot_maze.py
python -m pytest all_experiments/utils -q
\end{lstlisting}

\begin{utptable}
\begin{tabularx}{\linewidth}{@{}>{\bfseries\leavevmode\color{UTPNavy}}m{0.24\linewidth}m{0.36\linewidth}X@{}}
\toprule
\rowcolor{UTPPaleBlue!92}
\utpth{Command} & \utpth{You should see} & \utpth{It proves} \\
\midrule
The import line & \texttt{Packages OK} & The environment is active and complete \\
\texttt{demo\_robot\_maze.py} & Last line: \texttt{Every number above matches the book.} & Your copy of the search code behaves as the book describes \\
\texttt{pytest} & \texttt{8 passed} & The eight search tests pass on your machine \\
\bottomrule
\end{tabularx}
\end{utptable}

Lab~6 also needs the Tesseract OCR engine, which is not a Python package. Its lab sheet gives the install command for each system. Everything else needs only \texttt{requirements.txt}.

\begin{troubleshoot}
\begin{description}
  \item[\texttt{Activate.ps1} cannot be loaded] Windows blocks scripts by default. Run \texttt{Set-ExecutionPolicy -Scope CurrentUser RemoteSigned} once in PowerShell, then activate again.
  \item[Conda not found on Windows] Use Anaconda Prompt, or initialize PowerShell there with \texttt{conda init powershell}.
  \item[A notebook cannot import a package] Jupyter is running from a different environment. Activate the course environment, restart Jupyter, and run \texttt{import sys; print(sys.executable)} in a cell to confirm which Python it uses.
  \item[\texttt{No module named utils}] Start Jupyter from \texttt{all\_experiments/}, or run the \texttt{.py} files from the course folder as shown on each lab sheet.
\end{description}
\end{troubleshoot}
